\subsection{** APS‑TSLCA-REFERENCES-001 **}\label{apstslca-references-001}

\subsection{** Three-Squared-Lattice Cognitive Architecture **}\label{three-squared-lattice-cognitive-architecture}

\subsection{** Symbiotic Universal Xessability Standards **}\label{symbiotic-universal-xessability-standards}

\subsection{** Aurphyx Primordial Standard **}\label{aurphyx-primordial-standard}

\subsection{** Aurphyx LLC **}\label{aurphyx-llc}

\subsection{** SAGES \textbar{} Proprietary \textbar{} Pro-Existence **}\label{sages-proprietary-pro-existence}

\subsection{** Accessibility = Xessability **}\label{accessibility-xessability}

\subsection{** Version 3.69 **}\label{version-3.69}

\subsection{References}\label{references}

The references section is structured to match the conventions of formal scientific papers, particularly those submitted to arXiv in physics, mathematics, and computer science. Because the Three‑Squared‑Lattice Cognitive Architecture synthesizes concepts from differential geometry, tensor field theory, gauge theory, topological computation, distributed systems, and semantic formalisms, the references are grouped by domain. Each group contains the canonical foundational works that underpin the mathematical and theoretical structures used throughout the whitepaper.

\begin{center}\rule{0.5\linewidth}{0.5pt}\end{center}

\subsubsection{Foundational Mathematics and Geometry}\label{foundational-mathematics-and-geometry}

\begin{itemize}
\tightlist
\item
  S. Kobayashi and K. Nomizu, \emph{Foundations of Differential Geometry}, Vols. I--II, Wiley-Interscience.
\item
  M. Spivak, \emph{A Comprehensive Introduction to Differential Geometry}, Publish or Perish.
\item
  J. M. Lee, \emph{Introduction to Smooth Manifolds}, Springer.
\item
  R. Abraham, J. Marsden, and T. Ratiu, \emph{Manifolds, Tensor Analysis, and Applications}, Springer.
\item
  M. Nakahara, \emph{Geometry, Topology and Physics}, CRC Press.
\item
  A. Hatcher, \emph{Algebraic Topology}, Cambridge University Press.
\end{itemize}

\begin{center}\rule{0.5\linewidth}{0.5pt}\end{center}

\subsubsection{Tensor Fields, Gauge Theory, and Field Dynamics}\label{tensor-fields-gauge-theory-and-field-dynamics}

\begin{itemize}
\tightlist
\item
  C. Itzykson and J.-B. Zuber, \emph{Quantum Field Theory}, McGraw-Hill.
\item
  M. Peskin and D. Schroeder, \emph{An Introduction to Quantum Field Theory}, Addison-Wesley.
\item
  S. Weinberg, \emph{The Quantum Theory of Fields}, Vols. I--III, Cambridge University Press.
\item
  T.-P. Cheng and L.-F. Li, \emph{Gauge Theory of Elementary Particle Physics}, Oxford University Press.
\item
  A. Zee, \emph{Quantum Field Theory in a Nutshell}, Princeton University Press.
\end{itemize}

\begin{center}\rule{0.5\linewidth}{0.5pt}\end{center}

\subsubsection{Topological and Quantum Computation}\label{topological-and-quantum-computation}

\begin{itemize}
\tightlist
\item
  M. Freedman, A. Kitaev, M. Larsen, and Z. Wang, ``Topological Quantum Computation,'' \emph{Bull. Amer. Math. Soc.}
\item
  A. Kitaev, ``Fault-tolerant quantum computation by anyons,'' \emph{Annals of Physics}.
\item
  J. Preskill, \emph{Lecture Notes on Quantum Computation}, Caltech.
\item
  S. Lloyd, ``Universal Quantum Simulators,'' \emph{Science}.
\end{itemize}

\begin{center}\rule{0.5\linewidth}{0.5pt}\end{center}

\subsubsection{Dynamical Systems, Geometric Flows, and Thermodynamics}\label{dynamical-systems-geometric-flows-and-thermodynamics}

\begin{itemize}
\tightlist
\item
  S. Smale, ``Differentiable Dynamical Systems,'' \emph{Bull. Amer. Math. Soc.}
\item
  R. Hamilton, ``Three-manifolds with positive Ricci curvature,'' \emph{J. Differential Geometry}.
\item
  G. Perelman, ``The entropy formula for the Ricci flow and its geometric applications,'' arXiv:math/0211159.
\item
  H. Callen, \emph{Thermodynamics and an Introduction to Thermostatistics}, Wiley.
\end{itemize}

\begin{center}\rule{0.5\linewidth}{0.5pt}\end{center}

\subsubsection{Distributed Systems, Graph Theory, and Sheaf Theory}\label{distributed-systems-graph-theory-and-sheaf-theory}

\begin{itemize}
\tightlist
\item
  N. Biggs, \emph{Algebraic Graph Theory}, Cambridge University Press.
\item
  R. Diestel, \emph{Graph Theory}, Springer.
\item
  R. Ghrist, ``Sheaf Theory for Sensor Fusion,'' \emph{arXiv:1403.5809}.
\item
  L. Lovász, \emph{Large Networks and Graph Limits}, AMS.
\end{itemize}

\begin{center}\rule{0.5\linewidth}{0.5pt}\end{center}

\subsubsection{Semantic Systems, Identity, and Information Theory}\label{semantic-systems-identity-and-information-theory}

\begin{itemize}
\tightlist
\item
  C. Shannon, ``A Mathematical Theory of Communication,'' \emph{Bell System Technical Journal}.
\item
  G. Chaitin, \emph{Algorithmic Information Theory}, Cambridge University Press.
\item
  J. Barwise and J. Perry, \emph{Situations and Attitudes}, MIT Press.
\item
  S. Feferman, \emph{In the Light of Logic}, Oxford University Press.
\end{itemize}

\begin{center}\rule{0.5\linewidth}{0.5pt}\end{center}

\subsubsection{Cognitive Science, Embodiment, and Perception}\label{cognitive-science-embodiment-and-perception}

\begin{itemize}
\tightlist
\item
  A. Clark, \emph{Surfing Uncertainty: Prediction, Action, and the Embodied Mind}, Oxford University Press.
\item
  R. Friston, ``The free-energy principle: a unified brain theory?'' \emph{Nature Reviews Neuroscience}.
\item
  E. Thompson, \emph{Mind in Life}, Harvard University Press.
\item
  V. Gallese and G. Lakoff, ``The Brain's Concepts,'' \emph{Cognitive Neuropsychology}.
\end{itemize}

\begin{center}\rule{0.5\linewidth}{0.5pt}\end{center}

\subsubsection{Photonic and Hybrid Computing Substrates}\label{photonic-and-hybrid-computing-substrates}

\begin{itemize}
\tightlist
\item
  D. Miller, ``Silicon Photonics: Meshing Optics with Applications,'' \emph{Nature Photonics}.
\item
  J. Shainline et al., ``Superconducting Optoelectronic Circuits for Neuromorphic Computing,'' \emph{Physical Review Applied}.
\item
  H. Mabuchi, ``Coherent-feedback quantum control with photonic circuits,'' \emph{Physical Review A}.
\end{itemize}

\begin{center}\rule{0.5\linewidth}{0.5pt}\end{center}

\subsubsection{Ethical and Governance Frameworks}\label{ethical-and-governance-frameworks}

\begin{itemize}
\tightlist
\item
  L. Floridi and J. Sanders, ``On the Morality of Artificial Agents,'' \emph{Minds and Machines}.
\item
  N. Bostrom, \emph{Superintelligence}, Oxford University Press.
\item
  IEEE Global Initiative on Ethics of Autonomous and Intelligent Systems, \emph{Ethically Aligned Design}.
\end{itemize}

\begin{center}\rule{0.5\linewidth}{0.5pt}\end{center}

\subsubsection{Aurphyx Ecosystem Internal References}\label{aurphyx-ecosystem-internal-references}

These internal documents are part of the Aurphyx civilization‑scale architecture and are referenced throughout the whitepaper:

\begin{itemize}
\tightlist
\item
  \emph{Fuxyez: A Universal Semantic Transmutation Engine} (Whitepaper)
\item
  \emph{AuraFS: A Topological Filesystem for Proprioceptive Computation}
\item
  \emph{SAGES: The 13‑Invariant Governance Field}
\item
  \emph{Audry ABAS v1.0: Bicameral Synthetic Cognition}
\item
  \emph{APS-TSLCA-SUXS-USIS-SIX-001: Somatic Intelligence aXis Standard}
\item
  \emph{APS-TSLCA-SUXS-USIS-SCX-001: Systemic Coherence aXis Standard}
\item
  \emph{APS-TSLCA-SUXS-USIS-ICX-001: Identity Coherence aXis Standard}
\item
  \emph{APS-TSLCA-SUXS-USIS-SUXS-IFO-001: SUXS Intelligence Fusion Operator Standard}
\end{itemize}

\begin{center}\rule{0.5\linewidth}{0.5pt}\end{center}

\subsubsection{Notes on Citation Style}\label{notes-on-citation-style}

The references are structured to support:

\begin{itemize}
\tightlist
\item
  differential geometry\\
\item
  tensor field theory\\
\item
  gauge theory\\
\item
  topological computation\\
\item
  distributed cognition\\
\item
  semantic systems\\
\item
  embodied cognition\\
\item
  photonic substrates\\
\item
  ethical governance
\end{itemize}

This ensures the whitepaper is academically rigorous and cross‑disciplinary, matching the scope of the architecture.
